\documentclass[11pt]{article}

\usepackage[margin=1in]{geometry}
\usepackage{amsmath,amssymb,amsthm,mathtools}
\usepackage{enumitem}
\usepackage{booktabs}
\usepackage{tabularx}
\usepackage{microtype}
\usepackage[hidelinks]{hyperref}
\usepackage{xurl}

\newtheorem{definition}{Definition}
\newtheorem{lemma}{Lemma}
\newtheorem{remark}{Remark}

\title{Challenge Answer Carrier Slots for Successor Maintenance:\\[0.5ex]
\large Exact Shipped Fields and Reusable Clauses from Answer Cards}
\author{}
\date{Unified archive note v0.04 (2026-03-21; archive v482)}

\begin{document}
\maketitle

\begin{abstract}
Challenge answer cards already solve the released audience/question answer object.
They bind one whole-question handoff to one shipped sentence and one inherited visible-provenance mode.
But a later note that wants to audit the answer as actually carried still has to say where that sentence lives.
This note gives that last carrier question one home.
It defines a compact \emph{challenge answer carrier slot}: one exact shipped field or reusable clause that carries one emitted answer card without re-owning the card's semantics or reuse policy.
\end{abstract}

\paragraph{Non-redundancy note.}
Synthesis~31 owns the downstream suffix, Synthesis~35 owns sentence locks and inline-provenance policy, Synthesis~47 owns emitted answer cards, Synthesis~49 owns the full leftward answer-audit walk, Synthesis~54 owns the answer-side terminal citation normal form, Synthesis~74 owns the paired terminal citation discipline, Synthesis~75 owns the paired cutover rule, and Synthesis~17 owns the concrete worked instantiation.
This note owns only the exact carrier slot where one emitted answer lives once a later note wants to audit the shipped field or reusable clause itself.

\paragraph{Scope.}
This is not a new route map, stop profile, coverage certificate, answer sheet, answer card, sentence lock, or release stage.
It is a narrow carrier-interface note about binding one already-owned answer card to one already-owned field or clause.

\paragraph{Imports.}
Downstream clause and quote discipline are imported from Synthesis~31.\cite{series:downstreamnormalforms}
Sentence-lock policy is imported from Synthesis~35 only through the already-owned answer card.\cite{series:sentencelocks}
Challenge answer cards are imported from Synthesis~47.\cite{series:challengeanswercards}
The worked instantiation is imported from Synthesis~17.\cite{series:workedexample}
The answer-side terminal citation normal form and its paired reuse/cutover discipline are imported from Synthesis~54 and Synthesis~74--75 as the default later-paper stop rule around this exact-location owner.\cite{series:challengeanswerreviewmenus,series:pairedterminalcitationdiscipline,series:pairedterminalcutoverrules}

\paragraph{Exports.}
This note exports (i) challenge answer carrier slots, (ii) an exact-carrier rule, (iii) a card-carrier factorization rule, and (iv) an optional closure-role import rule saying that later notes may cite one exact shipped field or reusable clause carrying one emitted answer only when that exact carried location is itself live.
Otherwise later terse papers should default to the answer-side terminal citation normal form in Synthesis~54 together with the paired terminal discipline and paired cutover rule in Synthesis~74--75. Full leftward answer-audit walks remain outside scope and belong to Synthesis~49.\cite{series:challengeanswerreviewmenus,series:pairedterminalcitationdiscipline,series:pairedterminalcutoverrules}

\section{Why carrier slots deserve their own note}
There is a real difference between knowing what the released answer object is and knowing where it actually lives.
The answer card already owns the audience/question answer object.
But a maintainer, referee, or release note often needs one more fact: which exact field or reusable clause should be inspected to see that answer on the shipped surface.
Without one small carrier object, later notes keep restating some variant of ``take this answer card, then look at this wrapper field'' or ``take that answer card, then find the reusable clause entry carrying the same checked sentence.''
That is not just stylistic repetition.
It is a maintained binding between one answer object and one carrier slot, so the series benefits from giving it one narrow home.

\begin{definition}[Challenge answer carrier slot]
Fix one concrete base$\to$successor comparison, one audience key $a$, and one downstream question key $q$.
A \emph{challenge answer carrier slot} is a tuple
\[
K(a,q)=(a,q,C,s,\kappa,\lambda,\chi,\tau,p,\Phi)
\]
containing the audience key $a$, question key $q$, one imported answer-card key $C$, one carrier artifact path $s$, one carrier kind $\kappa\in\{\texttt{field},\texttt{clause}\}$, one carrier locator $\lambda$ inside that artifact, one carrier channel $\chi$, one exact carried surface text $\tau$, one imported inline-provenance profile key $p$, and the visible-provenance field list $\Phi$ inherited from the imported answer card.
\end{definition}

\begin{definition}[Exact-carrier rule]
A challenge answer carrier slot satisfies the \emph{exact-carrier rule} if the carried surface text $\tau$ agrees exactly with the text named by the declared carrier artifact and locator.
So the slot points to where the answer is actually carried; it does not paraphrase that answer a second time.
\end{definition}

\begin{definition}[Card-carrier factorization rule]
A challenge answer carrier slot satisfies the \emph{card-carrier factorization rule} if it references exactly one imported answer card and exactly one already-owned field or clause as the carrier, rather than restating piece order, sentence-lock guards, or visible-provenance policy inline.
The answer card continues to own the released answer object; the carrier slot only owns where that answer lives.
\end{definition}

\begin{definition}[Closure-role import rule]
If a challenge answer carrier slot declares a closure role, then it satisfies the \emph{closure-role import rule} if that role is already satisfied in the companion closure ledger for the same shipped artifact.
If no closure role is declared, then the slot is only a carrier binding and makes no new closure claim.
\end{definition}

\begin{lemma}[Carrier slots preserve answer identity while exposing exact location]
Assume every carrier slot satisfies the exact-carrier rule, the card-carrier factorization rule, and the closure-role import rule whenever a closure role is declared.
Then a later note may cite one carrier slot as the exact location of one released audience/question answer without changing the answer identity already owned by the imported answer card.
\end{lemma}

\begin{proof}
The exact-carrier rule prevents the slot from silently rewording the answer while claiming to point at its shipped location.
The card-carrier factorization rule keeps ownership separated: the imported answer card still owns the audience/question answer object, while the carrier slot owns only the field or clause where that answer appears.
The closure-role import rule prevents a slot from inventing a publication-completeness claim when it names a shipped wrapper role.
So the carrier slot adds only exact location, not new semantics or new reuse policy.\qedhere
\end{proof}

\begin{table}[t]
\centering
\small
\begin{tabularx}{\linewidth}{@{}l l X@{}}
\toprule
Layer & Canonical object & Question answered \\
\midrule
Released answer object & answer card & What audience/question answer is released, with which inherited sentence policy? \\
Exact carrier location & carrier slot & Which field or reusable clause actually carries that answer? \\
\bottomrule
\end{tabularx}
\caption{Carrier slots sit one level above answer cards: they expose exact location, not new answer semantics.}
\label{tab:answercarrierslots}
\end{table}

\section{Worked example pointer}
The maintained worked bundle now ships \path{example_successor_challenge_answer_carrier_slots.json}.\cite{series:workedexample}
It records four carrier slots for the canned Case-B successor.
The release-note public-status answer is carried by the shipped wrapper field \path{example_successor_lineage_notice.json:status_summary}, while the referee public-status answer is carried by the reusable clause \path{public_status_sentence} inside \path{example_successor_clause_pack.json}.
The release-note and auditor compact-diff answers are both carried by the reusable clause \path{compact_diff_clause} in that same clause pack.
Concrete keys already in the ledger include \path{release_note_public_status_sentence_answer_carrier_slot} and \path{auditor_compact_diff_summary_answer_carrier_slot}.
Each slot binds one answer card to one exact carrier artifact and locator, keeps the inherited visible-provenance profile explicit, and optionally imports a satisfied closure role when the carrier is itself a shipped wrapper field.

\paragraph{Why this helps the series.}
The stack is now cleaner at four distinct scales.
A capsule handles one semantic piece.
A sheet handles one whole audience/question handoff.
A card handles one released audience/question answer object.
And a carrier slot handles the exact field or reusable clause where that released answer actually lives.


\paragraph{A minimal public answer-carrier-slot card.}
\begin{table}[t]
\centering
\small
\begin{tabularx}{\linewidth}{@{}lX@{}}
\toprule
Field & Minimal public answer \\
\midrule
Release-note exact-location anchor & \nolinkurl{release_note_public_status_sentence_answer_carrier_slot} keeps answer card \nolinkurl{release_note_public_status_sentence_answer_card}, carrier artifact \nolinkurl{example_successor_lineage_notice.json}, carrier kind \texttt{field}, locator \nolinkurl{status_summary}, channel \texttt{public_wrapper}, default inline profile \texttt{floor}, visible provenance fields \nolinkurl{base_release_id}/\nolinkurl{successor_release_id}/\nolinkurl{compare_report_id}, and closure role \nolinkurl{fresh_successor_lineage_notice} \\
Auditor exact-location anchor & \nolinkurl{auditor_compact_diff_summary_answer_carrier_slot} keeps answer card \nolinkurl{auditor_compact_diff_summary_answer_card}, carrier artifact \nolinkurl{example_successor_clause_pack.json}, carrier kind \texttt{clause}, locator \nolinkurl{compact_diff_clause}, channel \texttt{reusable_clause}, and the same default visible-provenance floor \\
Selection rule & Stop here only while the live issue is the exact shipped field or reusable clause carrying one already-maintained released answer, not the narrower emitted-answer identity or the wider full leftward audit walk \\
Left/right boundary & Reopen Synthesis~47 when the live issue is the released answer object itself; widen to Synthesis~49 only when the live issue has become the full reader walk back through sheet, capsules, branches, and narrow terminal owners \\
Paired-terminal boundary & Once the checked answer and its exact location are no longer the moving part, later terse papers should usually stop at Synthesis~54 together with Synthesis~74--75 rather than reopen this carrier note merely to restate a now-fixed exact-location choice \\
\bottomrule
\end{tabularx}
\caption{A minimal public answer-carrier-slot card. This is the smallest answer-side public answer later terse papers can quote directly when the live issue is the exact shipped field or reusable clause carrying one already-maintained emitted answer.}
\label{tab:minimal-public-answer-carrier-slot-card}
\end{table}

This card is intentionally non-decorative: it pins the actual worked carrier-slot keys, their imported answer-card bindings, their exact artifact/locator identities, and the exact-location stop rule without silently re-owning emitted-answer identity, whole-question assembly, full audit, inline-budget policy, or paired-terminal cutover policy.\cite{series:workedexample}

\paragraph{A forced public answer-carrier-slot matrix.}
\begin{table}[t]
\centering
\small
\begin{tabularx}{\linewidth}{@{}lX@{}}
\toprule
Field & Forced public answer \\
\midrule
Release-note slot floor & \nolinkurl{release_note_public_status_sentence_answer_carrier_slot} still forces the exact location \nolinkurl{example_successor_lineage_notice.json:status_summary} only while the same emitted answer card \nolinkurl{release_note_public_status_sentence_answer_card}, the same carrier kind \texttt{field}, the same channel \texttt{public_wrapper}, the same visible-provenance field list, and the same declared closure role remain bound together \\
Auditor slot floor & \nolinkurl{auditor_compact_diff_summary_answer_carrier_slot} still forces the exact location \nolinkurl{example_successor_clause_pack.json:compact_diff_clause} only while the same emitted answer card \nolinkurl{auditor_compact_diff_summary_answer_card}, the same carrier kind \texttt{clause}, the same channel \texttt{reusable_clause}, and the same visible-provenance floor remain bound together \\
Imported-owner floor & Those shortcuts remain honest only while the same audience/question pair and the same imported answer-card key still certify the same emitted sentence lineage; the carrier slot does not get to keep the same exact-location citation after the answer card itself has changed \\
Locality floor & The carrier slot still only owns exact carried location; if a later paper needs the full leftward walk to capsules / branches / terminal owners, or needs new inline-budget / request-packet / reveal-order policy, this matrix has already widened out of the carrier-slot owner \\
Staleness trigger floor & This matrix goes stale as soon as the carrier artifact changes, the locator changes, the carrier kind or channel changes, the emitted answer card changes, the visible-provenance field list drifts, or a declared closure role no longer matches the companion closure ledger \\
Escalation pointer & Stop at Synthesis~47 when emitted-answer identity is still live; stop here only while one maintained slot still forces the same exact field or reusable clause; widen to Synthesis~49 only when the real issue becomes the full leftward reader walk rather than exact location alone \\
\bottomrule
\end{tabularx}
\caption{A forced public answer-carrier-slot matrix. This is the smallest answer-side public answer later terse papers can quote directly when the live issue is whether one maintained carrier slot still forces the same exact shipped location rather than merely which carrier slot exists in the abstract.}
\label{tab:forced-public-answer-carrier-slot-matrix}
\end{table}

This matrix is intentionally non-decorative: it pins the actual worked carrier-slot keys that still force one exact shipped field or reusable clause, together with the exact conditions under which that forcing remains honest or goes stale, without silently re-owning emitted-answer identity, full leftward audit, or later answer-side publication / review / terminal-cutover layers.\cite{series:workedexample}

\paragraph{One tiny answer-carrier-slot drill.}
For the worked release-note public-status answer, \nolinkurl{release_note_public_status_sentence_answer_carrier_slot} still forces \nolinkurl{example_successor_lineage_notice.json:status_summary} while the same emitted answer card, wrapper-field channel, visible provenance fields, and imported closure-role agreement remain intact. If the live question changes from ``where exactly is the answer carried?'' to ``walk every semantic piece leftward,'' the honest stop immediately widens to Synthesis~49 because the carrier slot no longer answers the whole question. If the live question stays on exact location but the emitted answer card, carrier locator, or wrapper-field closure role changes, the matrix has gone stale and a later terse paper must reopen the exact carrier-slot owner rather than keep quoting the old field path. Only once the exact carried location is fixed and the remaining question becomes the answer-side maintained handoff or paired terminal cutover should the paper leave this note for Synthesis~54 or Synthesis~74--75.\cite{series:workedexample}

\begin{thebibliography}{99}\small

\bibitem{series:downstreamnormalforms}
Anonymous.
\newblock Downstream Normal Forms for Successor Maintenance: Summary, Support, Citation, Quote, and Clause Discipline.
\newblock Series synthesis note (Synthesis~31), 2026. (This archive.)

\bibitem{series:sentencelocks}
Anonymous.
\newblock Sentence Locks and Inline Provenance for Successor Maintenance: Pair-Anchored Reuse, Recheck Guards, and Visible Capsules.
\newblock Series synthesis note (Synthesis~35), 2026. (This archive.)

\bibitem{series:challengeanswercards}
Anonymous.
\newblock Challenge Answer Cards for Successor Maintenance: Audience-Question Released Answers from Answer Sheets and Sentence Locks.
\newblock Series synthesis note (Synthesis~47), 2026. (This archive.)

\bibitem{series:challengeansweraudittrails}
Anonymous.

\newblock Challenge Answer Audit Trails for Successor Maintenance: Leftward Reader Walks from Carrier Slots to Narrow Terminals.

\newblock Series synthesis note (Synthesis~49), 2026. (This archive.)

\bibitem{series:challengeanswerreviewmenus}
Anonymous.

\newblock Challenge Answer Review Menus for Successor Maintenance: Smallest-Sufficient Referee Handoffs from Dockets, Profiles, Packets, and Ladders.

\newblock Series synthesis note (Synthesis~54), 2026. (This archive.)

\bibitem{series:pairedterminalcitationdiscipline}
Anonymous.

\newblock Paired Terminal Citation Discipline for Review Spines: Stable Checked-Answer Stops and Adjacent Request-Side Capstones.

\newblock Series synthesis note (Synthesis~74), 2026. (This archive.)

\bibitem{series:pairedterminalcutoverrules}
Anonymous.

\newblock Paired Terminal Cutover Rules for Review Spines: When to Stop, Widen, or Fail Closed.

\newblock Series synthesis note (Synthesis~75), 2026. (This archive.)

\bibitem{series:workedexample}
Anonymous.

\newblock Worked Example: Receipt Interlock for an Anonymous-DHT Reader-Privacy Claim.

\newblock Series synthesis note (Synthesis~17), 2026. (This archive.)

\end{thebibliography}
\end{document}
